\chapter{Pitted Keratolysis}
\index{Pitted Keratolysis}

\medskip
\noindent\textit{This chapter is for general education. A persistent foot odour or skin change should be assessed, especially with diabetes or reduced immunity.}

\section*{What it is}
Pitted keratolysis is a superficial bacterial infection of the outer skin, usually on pressure-bearing areas of the soles. It causes small pits, soft or macerated skin, and often a strong odour. Warm, wet conditions, sweaty feet, occlusive footwear, and prolonged standing or activity can contribute.

\section*{Diagnosis and treatment}
Diagnosis is often clinical. A clinician may test skin scrapings or consider other causes such as athlete’s foot or plantar warts. Treatment may include keeping feet dry, changing socks, breathable footwear, antiperspirant, antiseptic measures, and a prescribed topical antibiotic. Do not use leftover antibiotics or share footwear.

\section*{Self-care}
Wash and dry the feet carefully, rotate footwear so it can dry, change absorbent socks regularly, and avoid prolonged dampness. Follow advice for excessive sweating. People with diabetes should check the feet daily and seek early advice for breaks in the skin.

\section*{When to Seek Medical Care}
Arrange review for persistent pits, odour, pain, or treatment failure. Seek urgent care for rapidly spreading redness, warmth, swelling, fever, severe pain, a deep wound, or a cold or discoloured foot.

\section*{Sources and further reading}
\begin{itemize}
\item DermNet. \textit{Pitted keratolysis}. \url{https://dermnetnz.org/topics/pitted-keratolysis}
\end{itemize}
